\subsection{State-based construction of the twelve cyclotomic fibres}
\label{nuclear:fibras-app-witt}
On the APP marks, multiplication \(a:r\mapsto4r\) in
\(\ZZ/9\ZZ\) satisfies \(a^3=1\). The sheets satisfy
\[
 S(ax,ay)=aS(x,y),\qquad P(ax,ay)=a^{-1}P(x,y),
\]
since \(16\equiv7\equiv4^{-1}\pmod9\). The ordered unit orbits
are \(\mathcal O_0=(1,4,7)\), \(\mathcal O_1=(2,8,5)\).
In each, the augmentation module
\(A(\mathcal O)=\ker(\ZZ[\mathcal O]\xrightarrow{\sum}\ZZ)\)
inherits the inner product making the \(e_r\) orthonormal. In the basis
\(b_1=e_r-e_{4r}\), \(b_2=e_{4r}-e_{7r}\), its matrices are
\[
 G_{A_2}=\begin{pmatrix}2&-1\\-1&2\end{pmatrix},\qquad
 C=\begin{pmatrix}0&-1\\1&-1\end{pmatrix},\qquad
 C^{\mathsf T}G_{A_2}C=G_{A_2}.
\]
The lattice form therefore originates in the orbits of the marks.
For the pairs \((1,2),(2,3),(3,1)\) and sheets
\(\mu\in\{\Sigma,\Pi\}\), we obtain
\[
 W_{\APP}=\bigoplus_{(i,j)}\bigoplus_\mu\bigoplus_{\varepsilon=0}^1
 A(\mathcal O_\varepsilon),\qquad
 \ell^\Sigma_{ij}(x)=x_i+x_j,\quad
 \ell^\Pi_{ij}(x)=x_ix_j\pmod9.
\]
Each summand preserves pair, sheet, and orbit. The state map is
\[
 \Psi:(\ZZ/9\ZZ)^3\longrightarrow W_{\APP},\qquad
 \Psi(x)_{ij,\mu,\varepsilon}=\beta_\varepsilon(\ell^\mu_{ij}(x)),
 \quad
 \beta_\varepsilon(r)=
 \begin{cases}e_r-e_{4r},&r\in\mathcal O_\varepsilon,\\0,&r\notin\mathcal O_\varepsilon.
 \end{cases}
\]
\begin{proposition}[Equivariance and rank of the state module]
\label{nuclear:psi-rango}
The action of \(C\) on the six additive fibres and of \(C^{-1}\) on
the six multiplicative fibres satisfies \(\Psi(ax)=\rho(a)\Psi(x)\).
The images span \(W_{\APP}\otimes\QQ\), of dimension twenty-four.
\end{proposition}
\begin{proof}
The sheet identities and \(\beta(4r)=C\beta(r)\) prove
equivariance. For rank, order the coordinates by the three
stated pairs, by \(\Sigma,\Pi\), by \(\mathcal O_0,\mathcal O_1\),
and by \(b_1,b_2\). In each orbit, \(\beta\) takes successive values
\((1,0),(0,1),(-1,-1)\). The matrix of the images of the following states,
read row by row, has determinant \(9\):
\[
\begin{array}{rrrr}
(0,0,1)&(0,0,2)&(0,0,4)&(0,0,5)\\
(0,1,0)&(0,1,1)&(0,1,2)&(0,1,3)\\
(0,1,4)&(0,1,5)&(0,1,6)&(0,2,0)\\
(0,2,2)&(0,2,3)&(0,4,0)&(0,5,0)\\
(1,0,1)&(1,0,2)&(1,0,4)&(1,0,5)\\
(1,1,0)&(1,2,0)&(1,4,0)&(1,5,0).
\end{array}
\]
Successive elimination, reducing each row by its predecessors and
normalizing its first nonzero coefficient, gives pivot columns
numbered from zero
\[
\begin{split}
&(8,10,9,11,0,12,14,16,13,15,17,2,\\
&\hspace{2em}18,19,1,3,20,22,21,23,4,6,5,7)
\end{split}
\]
and pivots
\[
(1,1,1,-1,1,1,1,1,1,-1,3,1,1,3,1,-1,1,1,1,-1,1,1,1,-1).
\]
Their product, multiplied by the sign of the column permutation,
is \(9\). The invertible minor establishes full rank.
\end{proof}

\subsection{Dodecaphasic alignment with the Witt positions}
\label{nuclear:alineamiento-witt}
The pair index is \(i\in\ZZ/3\ZZ\), the sheet index is
\(\mu\in\{0,1\}\), and the orbit index is \(\varepsilon\in\{0,1\}\).
With the dodecaphasic origin fixed, \(\lambda(i,\mu,\varepsilon)
=4i+2\mu+\varepsilon\pmod{12}\) is bijective: the four values
for each \(i\) form a block, and the three blocks are disjoint.
Let \(R=\left(\begin{smallmatrix}0&1\\1&0\end{smallmatrix}\right)\).
Multiplication proves \(RCR=C^{-1}\) and
\(R^{\mathsf T}G_{A_2}R=G_{A_2}\). For the trivializations
\(\iota_b\) and frames \(P_b\) transported jointly with
Witt incidence, set
\[
g_b=\iota_b^{-1}P_bCP_b^{-1}\iota_b,\qquad
T_{i\mu\varepsilon}=\iota_{\lambda(i,\mu,\varepsilon)}^{-1}
P_{\lambda(i,\mu,\varepsilon)}R^\mu.
\]
Then
\[
g_{\lambda(i,\mu,\varepsilon)}T_{i\mu\varepsilon}
=T_{i\mu\varepsilon}C^{(-1)^\mu}.
\]
The direct sum is a \(C_3\)-equivariant isomorphism between the twelve
state fibres and the twelve lattice positions. It preserves pair,
sheet, orbit, and form under simultaneous transport of origin,
order, and orientation. These data persist in the lattice gluing
and passage to the neighbour constructed below.
